(1) 设链的水平张力分量为常数 \(H>0\)。弧长元 \(ds=\sqrt{1+y'^2}\,dx\)，竖直平衡为 \(d(Hy')/dx=\lambda g\sqrt{1+y'^2}\)。令 \(a=H/(\lambda g)\)，积分得到 \(\operatorname{arsinh}y'=x/a\)，其中对称最低点取在 \(x=0\)。因此 \(y'=\sinh(x/a)\)，再由两端 \(y(\pm\ell/2)=0\) 得
\[
\boxed{y(x)=a[\cosh(x/a)-\cosh(\ell/(2a))]},\qquad H=\lambda ga.
\]
该曲线严格来自张力平衡而非先假设抛物线近似。

(2) 由 \(ds=\cosh(x/a)dx\)，全链长
\[
\boxed{S=\int_{-\ell/2}^{\ell/2}\cosh(x/a)dx=2a\sinh(\ell/(2a))}.
\]
令 \(u=\ell/(2a)>0\)，则 \(S/\ell=\sinh u/u\)。函数导数 \((u\cosh u-\sinh u)/u^2>0\)，因为分子在0为零、导数为 \(u\sinh u>0\)。且 \(u\to0\) 时函数趋1，\(u\to\infty\) 时趋无穷，故每个 \(S>\ell\) 唯一确定 \(u\)、\(a\) 与 \(H\)。

(3) 最低点垂度
\[
d=-y(0)=a[\cosh(\ell/(2a))-1].
\]
张力模长 \(T(x)=H\sqrt{1+y'^2}=H\cosh(x/a)\)，最低点最小张力为 \(H\)。支点竖直张力分量的大小为 \(H\sinh(\ell/(2a))=\lambda g a\sinh(\ell/(2a))=\lambda gS/2\)，两个支点合计承担 \(\lambda gS\)，与全链重量严格相等；支点水平分量互相抵消。

(4) 以支点高度为零势，链势能
\[
U=\lambda g\int_{-\ell/2}^{\ell/2}y(x)\sqrt{1+y'^2}\,dx
=\boxed{\lambda ga^2[u-\sinh u\cosh u]},\quad u=\ell/(2a).
\]
该值为负，因链位于支点下方。浅垂度 \(u\ll1\) 时
\(S/\ell=\sinh u/u=1+u^2/6+O(u^4)\)，
\(d=a(\cosh u-1)=\ell u/4+O(\ell u^3)\)，
\(H=\lambda g\ell/(2u)\)。因此
\[
d\simeq\frac{\ell}{4}\sqrt{6(S/\ell-1)},\qquad
H\simeq\frac{\lambda g\ell}{2\sqrt{6(S/\ell-1)}}.
\]
当 \(S\downarrow\ell\) 时张力发散，显示不可伸长柔链的理想化极限。
